\section{Conclusion and Future Work}
\label{sec:conclusion}

In this work, we deconstructed the capacity-generalization trade-off in dynamic weight-space model merging and challenged the prevailing practice of using high-capacity, unshared layer-wise routing networks. We mathematically formalized and empirically proved the phenomenon of \textbf{coefficient ruggedness}, where fully independent routing weights across sequential layers introduce excessive degrees of freedom, leading to high layer-to-layer discrepancy, optimization variance, and cascading representation drift. 

To resolve this issue, we introduced the \textbf{Block-wise Weight-Sharing Router} (\textbf{BWS-Router}), which groups the $L$ layers of the network into $G = L / M$ uniform blocks and shares routing parameters within each block. Under an exhaustive validation sweep of over $1,280$ experiments across five independent seeds inside our task-conflict sandbox, we demonstrated that BWS-Router ($M=3$) achieves an outstanding dynamic Joint Mean accuracy of \textbf{79.56 $\pm$ 1.13\%} (and climbing to \textbf{79.63 $\pm$ 1.18\%} under optimal learning rates), compressing the parameter footprint of unshared baselines by 66.7\% while massively outperforming static uniform merging (which collapses to \textbf{23.56 $\pm$ 2.91\%}). Furthermore, we validated BWS-Router under a **true physical sequential weight-space model-merging framework**, where our block-shared router ($M=3$) achieves \textbf{45.26 $\pm$ 10.11\%} Joint Mean accuracy, outperforming static uniform merging (\textbf{17.88 $\pm$ 3.78\%}), and significantly boosting heterogeneous mixed-batch accuracy by \textbf{+10.93\%} absolute over the unshared baseline (achieving \textbf{43.20 $\pm$ 22.49\%} vs. \textbf{32.27 $\pm$ 21.28\%}). This confirms that parameter sharing within blocks acts as a powerful structural regularizer that stabilizes optimization and prevents representation drift inside physical weight-space ensembling systems. 

Additionally, we demonstrated that BWS-Router possesses remarkable resiliency under task-heterogeneous batch shifts, maintaining a stable accuracy of \textbf{79.30 $\pm$ 1.88\%} when evaluated sample-by-sample, and that scaling learning rates and regularizing routing weights prevents OOD SVHN performance collapse. Finally, we exposed the overparameterization and complexity of wave-based routing (QWS-Merge), establishing regularized block weight-sharing as a far superior alternative.

\paragraph{Bridge to Physical Model Merging: Implementation Recipe for Deep ViTs}
To facilitate the transition from our controlled sandbox simulation to real weight-space model merging, we propose a concrete implementation recipe for BWS-Router on physical deep Vision Transformers (ViTs, e.g., ViT-B/16). 
First, given $K$ expert checkpoints fine-tuned from a shared pre-trained initialization, we construct task vectors $\Delta W_k^{(l)} = W_k^{(l)} - W_{pre}^{(l)}$ for every layer $l \in \{1, \dots, L\}$. 
Second, we partition the $L = 12$ Transformer blocks into $G = 4$ uniform blocks of size $M = 3$. For each block group $g \in \{1, \dots, G\}$, we learn a single, shared router parameterized by $W_{group}^{(g)}$. 
Third, unlike our virtual-layer sandbox, representation propagation in a physical ViT is sequential: the input representation $x^{(l)}$ fed into the first layer of block group $g$ is used to predict the gating coefficients $\alpha_{g, k}(x^{(l)}) = \text{Sigmoid}(W_{group}^{(g)} \cdot \text{PCA}^{(g)}(x^{(l)}))$. To maximize sequential alignment and capture the shifting representation manifold across network depths, the unsupervised block-specific PCA projectors $\text{PCA}^{(g)}$ are fit sequentially during the calibration forward pass (fitting each block group's projector sequentially on the intermediate calibration hidden features outputted by the preceding block group $g-1$ under uniform weight blending).
Fourth, the blended weights for all layers $l$ within group $g$ are computed at runtime as $W^{(l)}_{merged}(x) = W_{pre}^{(l)} + \sum_{k=1}^K \alpha_{g, k}(x^{(l)}) \cdot \Delta W_k^{(l)}$.
Finally, the router parameters are optimized using a tiny calibration set of downstream task samples. Because sequential backpropagation through 12 unshared layers suffers from high gradient variance and exploding/vanishing gradients under data-scarce adaptation, sharing routing weights within blocks structurally restricts the optimization space, smoothing the loss landscape and ensuring stable convergence. Furthermore, to stabilize compounding representation drift and reduce seed-wise performance variance under sequential propagation, we explicitly advise practitioners to employ \textbf{sequential smoothing regularization} ($\mathcal{L}_{\text{smooth}}$) rather than \textbf{residual routing links}. As evaluated in our sweeps (see Appendix~\ref{app:smoothing}), residual routing links successfully reduce seed variance but severely degrade absolute accuracy by forcing routing coefficients toward a task-agnostic static average. In contrast, sequential smoothing regularization stabilizes deep optimization trajectories and reduces seed-wise standard deviation from \textbf{21.28\%} down to \textbf{13.41\%} while fully preserving runtime routing expressivity and dynamic accuracy ceilings.

\paragraph{Limitations, Architectural Scalability, and Future Directions}
While BWS-Router provides a highly efficient and stable dynamic merging framework, there are several limitations and architectural design questions that represent critical pathways for future research and deployment. To satisfy strict page limits, a detailed discussion on bridging BWS-Router to physical checkpoints, scaling to larger expert counts ($K \ge 10$), choosing Softmax vs. Sigmoid gating rules, scaling PCA projectors across depths, mitigating representation drift, and dynamic layer grouping is deferred to Appendix~\ref{app:limitations_scalability}.
